\chapter{Conclusion and Future Work}
\label{ch:conclusion}

This thesis has presented a comprehensive investigation of multi-paradigm deep learning approaches for financial fraud detection, systematically comparing Graph Neural Network (GNN), Spiking Neural Network (SNN), and Autoencoder (AE) ensemble pipelines on the Bank Account Fraud (BAF) dataset \cite{jesus2022bank,feedzai2022baf}. The research was motivated by the observation that no single modeling paradigm is universally optimal for fraud detection \cite{singh2024comprehensive}, and that the relative strengths of different architectures remain poorly understood due to the lack of controlled cross-paradigm comparisons in the existing literature \cite{pourhabibi2022graph}. This chapter summarizes the contributions, offers critical reflection on the findings, and outlines recommendations for future research.

\section{Summary of Contributions}

This thesis makes the following eight contributions to the field of deep learning-based financial fraud detection:

\begin{enumerate}[label=\arabic*.]
\item \textbf{FraudGraphBuilder Framework:} We designed and implemented a robust graph construction pipeline that transforms flat tabular fraud data into a heterogeneous graph with 17~node types, over 1~million nodes, and 14.5~million edges. The framework incorporates domain-informed feature bucketing strategies, graceful NaN handling, and edge symmetrization for bidirectional message passing. This framework serves as the foundation for the entire GNN pipeline and could be adapted to other domains where tabular data contains relational patterns worth exploiting through graph representations.

\item \textbf{Multi-Paradigm Architecture Comparison:} We implemented and systematically compared five deep learning architectures across three paradigms---GraphSAGE and QGNN (spatial/quantum GNN), Split GNN (spectral GNN), SNN with rate coding and ensemble distillation (neuromorphic), and VAE-Stacked and DAE-Stacked autoencoders (unsupervised/stacking)---on the same fraud detection dataset. This represents the most comprehensive cross-paradigm comparison for financial fraud detection to date, revealing that the SNN Variant~III achieves the highest test ROC AUC (0.9260), followed by the VAE-Stacked (0.8868) and GraphSAGE (0.8859).

\item \textbf{Comprehensive Hyperparameter Sensitivity Analysis:} We conducted a 45-experiment grid search for the Split GNN, revealing that simpler spectral filtering ($C=1$) consistently outperforms complex alternatives, that incorporating same-class edge information ($K=0$) is essential, and that moderate edge supervision ($\gamma_{\text{edge}} = 0.4$) provides the most robust performance. The 4.04 percentage point AUC gap between the best and worst configurations demonstrates that hyperparameter tuning is not optional but essential for spectral GNN architectures.

\item \textbf{SNN Pipeline with Rate Coding and Ensemble Methods:} We developed a spiking neural network pipeline for tabular fraud detection that employs Bernoulli (Poisson-like) rate coding, leaky integrate-and-fire (LIF) neurons with surrogate gradient training, and an ensemble distillation strategy (Variant~III) that combines multiple threshold configurations. This pipeline achieves the best overall performance (AUC 0.9260) and the highest robustness to distributional shift (CV of 5.53\%), demonstrating the potential of neuromorphic computing principles for financial applications \cite{nasif2026reinforcement}.

\item \textbf{Dual-Autoencoder Stacking Framework:} We proposed a novel stacking framework that combines a Variational Autoencoder (VAE) \cite{kingma2014auto} and a Denoising Autoencoder (DAE) \cite{vincent2010stacked} as base anomaly detectors with a gradient-boosted meta-learner \cite{chen2016xgboost,ke2017lightgbm}. The VAE-Stacked configuration achieves an AUC of 0.8868, demonstrating that unsupervised anomaly detection combined with supervised stacking can nearly match the performance of purely supervised approaches. The target manifold isolation strategy effectively sidesteps the class imbalance problem during representation learning.

\item \textbf{Cross-Version Fairness Evaluation:} We evaluated model robustness across six data distributions, demonstrating substantial performance variability (GNN AUC range: 0.7587--0.9454; SNN AUC range: 0.8101--0.9587; AE AUC range: 0.7893--0.9421) and identifying the critical risk of distributional shift in production environments. The generalization gap metric proved to be an effective early warning indicator. The SNN pipeline exhibited the highest stability (CV of 5.53\% vs.\ 7.21\% for GraphSAGE), suggesting that temporal spike dynamics provide natural resilience to distributional perturbations.

\item \textbf{Resampling Strategy Evaluation:} We demonstrated that inverse-frequency class weighting outperforms ADASYN-based resampling \cite{he2008adasyn,chawla2002smote} for graph-based fraud detection, providing better-calibrated probability outputs (threshold 0.7621 vs.\ 0.9944) and superior overall performance (AUC 0.8859 vs.\ 0.8748). This finding extends across pipelines, as distribution-modifying strategies \cite{farahat2025objective,dong2025enhanced} are particularly harmful for graph-based models where neighborhood structure is critical, and unreliable for anomaly-based models where synthetic minority samples may not faithfully represent the fraud manifold. The resampling paradox identified by Dong \cite{dong2025enhanced}---where synthetic oversampling within autoencoder training loops causes severe probability miscalibration---confirms that loss-based and architecture-based approaches are preferable across all paradigms.

\item \textbf{Cross-Pipeline Comparison and Deployment Recommendations:} We provided the first systematic cross-pipeline comparison of GNN, SNN, and AE approaches for financial fraud detection, offering actionable recommendations for architecture selection, monitoring, class imbalance handling, and fairness-aware deployment. The analysis reveals that the three pipelines capture complementary fraud signals---relational patterns (GNN), temporal dynamics (SNN), and distributional anomalies (AE)---suggesting that an ensemble of all three paradigms could yield further performance improvements.
\end{enumerate}

\section{Critical Reflection}

Several aspects of this research warrant critical reflection to provide a balanced assessment of the contributions and their implications:

\textbf{Performance vs.\ Complexity Trade-off:} The SNN Variant~III's superior performance comes at the cost of increased architectural complexity. The rate coding transformation, temporal simulation across multiple time steps, and ensemble distillation require careful implementation and hyperparameter selection. In contrast, the GraphSAGE model achieves competitive performance with a much simpler architecture and faster training time. For organizations without deep learning expertise or the computational resources to train SNN models, the GraphSAGE or AE-Stacked pipelines---or indeed, well-tuned gradient-boosted tree models \cite{grinsztajn2022tree}---may represent more pragmatic choices despite their lower peak performance.

\textbf{SNN Calibration Challenges:} Despite achieving the highest AUC, the SNN pipeline presents calibration challenges that are not fully addressed in this work. The surrogate gradient training introduces an approximation to the true gradient that can lead to miscalibrated probability outputs, where the predicted probabilities do not accurately reflect the true likelihood of fraud. This is particularly problematic in financial applications where probability estimates inform risk-based decision-making and regulatory reporting. Future work should incorporate calibration techniques such as temperature scaling or Platt scaling specifically adapted for SNN outputs.

\textbf{Autoencoder Standalone Insufficiency:} The autoencoder pipeline's reliance on the stacking meta-learner raises questions about the standalone value of the anomaly detection approach. As individual anomaly detectors, the VAE and DAE produce AUC values of approximately 0.72--0.78, which are insufficient for production deployment. This aligns with the broader literature: Dong \cite{dong2025enhanced} reported a standalone DAE AUC of only 0.5210 on the BAF dataset, while Obushyni et al. \cite{obushyni2025vae} found that standard autoencoders achieve AUC of approximately 0.55, compared to 0.8954 for VAEs with probabilistic scoring. The stacking framework effectively rescues the autoencoder representations by combining them with supervised learning, but this diminishes the appeal of the unsupervised approach for cold-start scenarios where labeled data is truly unavailable. The unsupervised anomaly score alone cannot replace a supervised classifier.

\textbf{Fairness Concerns Across Pipelines:} While the cross-version evaluation demonstrates robustness to distributional shift, the subgroup fairness analysis reveals persistent age-based TPR disparities of up to 13.6\% across all pipelines. The SNN pipeline shows the smallest disparity (12.1\%), but this still represents a significant inequity that could have legal and ethical implications in regulated financial services. The consistent nature of this disparity across all three paradigms suggests that it may be rooted in the underlying data rather than in model-specific biases \cite{chouldechova2017fair,kozodoi2022fairness}, which would require data-level interventions rather than algorithmic fixes alone.

\textbf{Evaluation Limitations:} The 5\% FPR operating point, while practical, represents a single point on the ROC curve. A more comprehensive evaluation would consider the full cost curve, incorporating the financial cost of fraud losses against the operational cost of false positive investigations. Different pipelines may have different optimal operating points depending on the cost structure of the deployment context. Additionally, the AUC metric, while widely used, can be misleading for highly imbalanced datasets where the precision-recall curve provides a more informative picture of model performance.

\textbf{Reproducibility and Computational Cost:} While all experiments are documented with fixed hyperparameters and random seeds, the computational cost of the full experimental suite across all three pipelines (estimated 400+ GPU-hours) may limit independent verification and rapid iteration. The Split GNN alone requires 45~configurations totaling approximately 19~hours of GPU time, and the SNN ensemble distillation adds significant overhead. Researchers and practitioners with limited computational budgets may need to prioritize specific pipelines or use more efficient hyperparameter search strategies such as Bayesian optimization.

\textbf{Dataset-Specific Findings:} All findings are derived from a single dataset suite (BAF). While the multi-version design of BAF provides some diversity, the specific characteristics of credit card application fraud may not generalize to other fraud domains. For instance, the graph-based approach's advantage in capturing relational patterns may be more pronounced in domains with denser entity relationships (e.g., money laundering networks), while the SNN's temporal processing may be more beneficial in domains with strong temporal dynamics (e.g., real-time transaction monitoring).

\section{Recommendations for Future Research}

Based on the findings and limitations of this thesis, we propose the following ten recommendations for future research:

\begin{enumerate}[label=\arabic*.]
\item \textbf{Temporal Graph Neural Networks:} Extending the current static graph representation to a dynamic temporal graph that captures the evolution of fraud patterns over time would address a fundamental limitation of the GNN pipeline. Temporal GNN architectures such as TGAT (Temporal Graph Attention) and EvolveGCN could provide additional signal from temporal dynamics, potentially closing the performance gap with the SNN pipeline. The temporal graph would capture not just who is connected to whom, but when those connections were established and how they change over time, providing a richer representation for detecting fraud campaigns that unfold over weeks or months.

\item \textbf{Real Quantum Hardware Evaluation:} As quantum hardware becomes more accessible, evaluating the QGNN architecture on actual quantum processors (e.g., IBM Quantum, Google Sycamore) would provide definitive evidence on the quantum advantage for fraud detection. The current simulation-based results may not reflect the true potential of quantum feature maps, which exploit genuine quantum entanglement and superposition that classical simulations can only approximate. The barren plateaus phenomenon may also manifest differently on real hardware, potentially requiring new optimization strategies such as layer-wise training or quantum natural gradients.

\item \textbf{Heterogeneous GNN Architectures:} The current GNN models treat all edge types uniformly despite the heterogeneous nature of the fraud detection graph. Heterogeneous GNN architectures such as Relational Graph Convolutional Networks (RGCN) and Heterogeneous Graph Transformers (HGT) that explicitly model edge type differences could better leverage the rich node type structure of the fraud detection graph. Different edge types (shared device, shared address, similar risk profile) carry different diagnostic values, and a model that can learn type-specific aggregation functions may achieve superior performance.

\item \textbf{Neuromorphic Hardware Deployment:} The SNN pipeline's high inference latency on standard GPU hardware ($\sim$50~ms per prediction) is a significant deployment barrier. Evaluating the SNN on neuromorphic hardware such as Intel Loihi or IBM TrueNorth could dramatically reduce inference latency and energy consumption, potentially making SNN-based fraud detection viable for real-time deployment. The event-driven nature of neuromorphic processors aligns naturally with the spike-based computation model, and early benchmarks suggest 100--1000$\times$ improvements in energy efficiency for SNN inference on neuromorphic hardware.

\item \textbf{Advanced Autoencoder Variants:} Exploring more sophisticated autoencoder architectures could improve the standalone anomaly detection performance. Adversarial autoencoders, which combine the variational framework with adversarial training to enforce a prior distribution on the latent space, may produce more discriminative latent representations. Normalizing flows and diffusion-based models offer alternative approaches to density estimation that could provide more principled anomaly scores. Additionally, conditional VAEs that incorporate protected attributes as conditioning variables could simultaneously improve detection performance and fairness.

\item \textbf{Explainability and Interpretability:} Developing explainability methods for all three pipelines is critical for building trust and facilitating adoption by fraud analysts. For the GNN pipeline, edge importance visualization and subgraph explanation methods (e.g., GNNExplainer, PGExplainer) could identify which relational patterns the model considers most indicative of fraud. For the SNN pipeline, spike timing analysis and neuron importance scoring could reveal which temporal patterns trigger fraud detection. For the AE pipeline, latent space visualization and reconstruction error attribution could help analysts understand why specific applications are flagged as anomalous.

\item \textbf{Fairness-Aware Training:} Incorporating fairness constraints directly into the training objective across all pipelines could reduce the demographic disparities observed in the subgroup fairness analysis \cite{kozodoi2022fairness}. Adversarial debiasing, where a secondary network tries to predict protected attributes from the model's internal representations while the primary network tries to prevent this, could learn representations that are invariant to age and employment status. Equalized odds regularization, which penalizes differences in TPR and FPR across demographic groups, could be added as a loss term during training. Pre-processing approaches such as fair representation learning could also be explored.

\item \textbf{Production Deployment and MLOps:} Developing a production-grade deployment pipeline with model serving, A/B testing, automated retraining, and performance monitoring would bridge the gap between research prototype and operational system. The cross-version evaluation suggests that model retraining frequency should be aligned with the observed rate of distributional shift, and that monitoring systems should track both overall metrics and subgroup-specific metrics to detect fairness drift. The ensemble of GNN, SNN, and AE predictions could be deployed as a composite model with configurable decision thresholds for different operational contexts.

\item \textbf{Cross-Domain Validation:} Evaluating the proposed multi-pipeline framework on fraud detection datasets from other domains---insurance claims, e-commerce transactions, cryptocurrency fraud, healthcare fraud---would assess the generalizability of the findings. The relative performance of the three pipelines may shift depending on domain-specific characteristics: graph-based approaches may dominate in domains with rich entity relationships (money laundering), SNNs may excel in domains with strong temporal dynamics (transaction monitoring), and autoencoders may be preferred in domains with limited labeled data (emerging fraud types). A cross-domain evaluation would provide more nuanced deployment guidance.

\item \textbf{Adversarial Robustness:} Evaluating the robustness of all three pipelines against adversarial attacks is critical for security-sensitive financial applications. Graph structure attacks (edge injection, edge deletion) could deceive the GNN pipeline by creating fake connections that mask fraud rings. Feature perturbation attacks could modify application attributes to evade detection by the SNN and AE pipelines. Evaluating adversarial robustness across pipelines would reveal which architectures are most vulnerable and inform the design of adversarial training strategies that improve resilience against sophisticated fraudsters who actively attempt to circumvent detection systems.
\end{enumerate}
